\documentclass[12pt]{article}

\usepackage[a4paper, margin=2.5cm]{geometry}
\usepackage{amsmath}
\usepackage{amssymb}
\usepackage{amsthm}
\usepackage[colorlinks=true, linkcolor=blue, citecolor=blue, urlcolor=blue]{hyperref}
\usepackage{booktabs}
\usepackage{natbib}

\newtheorem{theorem}{Theorem}
\newtheorem{proposition}[theorem]{Proposition}
\newtheorem{corollary}[theorem]{Corollary}
\theoremstyle{definition}
\newtheorem{definition}[theorem]{Definition}
\theoremstyle{remark}
\newtheorem{remark}[theorem]{Remark}

\newcommand{\Vcal}{\mathcal{V}}

\title{Hubble Constant and $\Omega_\Lambda$ from the Cascade\\[6pt]
\large Cosmological Parameters as Projections of the $E_8$ Closure Depth}
\author{%
  Lee Smart\\[4pt]
  \textit{Institute of Vibrational Field Dynamics}\\[2pt]
  \texttt{contact@vibrationalfielddynamics.org}\\[2pt]
  \texttt{@vfd\_org}%
}
\date{April 2026}

\begin{document}

\maketitle

\noindent\textbf{Status:} Preprint (not peer-reviewed).

\begin{abstract}
The $\Lambda$CDM cosmological model is correct but incomplete: its six parameters ($H_0$, $\Omega_\Lambda$, $\Omega_m$, $n_s$, $A_s$, $\tau$) are fit to CMB + large-scale-structure + supernova data rather than derived. Two of these --- the Hubble constant $H_0$ and the dark-energy density parameter $\Omega_\Lambda$ --- have known ``tensions'' between early-universe (Planck CMB) and late-universe (SH0ES distance ladder) measurements, indicating that the standard framework does not internally resolve its own parameters.

This paper derives $H_0$ and $\Omega_\Lambda$ from the same cascade foundations that produced the cosmological constant $\Lambda \cdot \ell_P^2 = 2\varphi^{-583}$ in Paper XXXVI~\citep{SmartXXXVI}. Under F1 ($r = \varphi$), F4 (closure depth $583$), F5 (factor 2), and F7 (rank-2 $\Lambda$ on $D_4$), the Hubble constant and $\Omega_\Lambda$ emerge as cascade-structural consequences, not fit parameters:

\begin{itemize}
\item $H_0 = m_P \cdot \varphi^{-291.5} \approx 68.83$ km/s/Mpc (derivation: \texttt{cascade-lambda.md §9.6}). Observed range: Planck 2018 gives $67.36 \pm 0.54$ km/s/Mpc; SH0ES 2022 gives $73.04 \pm 1.04$ km/s/Mpc. The cascade value sits centrally inside the Hubble-tension window.
\item $\Omega_\Lambda \approx 2/3$ from dual-600-cell energy-balance ratio. Observed value $\Omega_\Lambda = 0.685 \pm 0.007$ (Planck 2018): $2.7\%$ agreement.
\item Invariant $H_0 \cdot \sqrt{\Omega_\Lambda}$: cascade prediction $56.9$, observed $56.4$ (Planck): $0.81\%$ agreement. This invariant is insensitive to the Hubble tension.
\item Cosmic age $t_0 = 13.81$ Gyr: $0.1\%$ agreement with Planck 2018.
\end{itemize}

VFD does not replace $\Lambda$CDM. It derives the cosmological parameters from the cascade structure already used to produce $\Lambda$, yielding a parameter-free cosmology consistent with observation at sub-percent precision on the invariant $H_0\sqrt{\Omega_\Lambda}$ and with the Hubble-tension centre on $H_0$ itself.
\end{abstract}

%==================================================================
\section{Introduction and Recovery Position}\label{sec:intro}
%==================================================================

\textbf{Recovery position.} This paper carries Recovery Items beyond R8 (cosmological constant) of the Recovery Manifest~\citep{RecoveryManifest}: specifically, the Hubble constant $H_0$, the dark-energy density $\Omega_\Lambda$, and the cosmic age $t_0$. Upstream: Paper XXXVI (foundations F1, F4, F5, F7), Paper XXXV (three-pillars synthesis). Downstream: Paper LI (inflation and dark matter), which extends the cascade to the early-universe and dark-matter sectors. Authoritative first-principles derivation: \texttt{papers/cascade-derivation/cascade-lambda.md} (1144 lines), specifically §9.6 ($H_0$), §9.7 ($\Omega_\Lambda$), §9.10 (falsifiable predictions including cosmic age).

The $\Lambda$CDM model has six free parameters fit to observation. This paper reduces two of them ($H_0$, $\Omega_\Lambda$) to consequences of cascade foundations F1--F7. Combined with Paper XXXVI's derivation of $\Lambda$ itself and Paper XXXV's three-pillars synthesis, the cascade delivers the full late-universe cosmological parameter set ($\Lambda$, $H_0$, $\Omega_\Lambda$, $t_0$) from two axioms plus classical mathematics.

%==================================================================
\section{The Standard Cosmological Parameters}\label{sec:standard}
%==================================================================

\textbf{Standard form.} $\Lambda$CDM describes the late-time homogeneous isotropic universe with Friedmann equation
\begin{equation}\label{eq:friedmann}
H^2(a) = H_0^2 \left[ \Omega_m a^{-3} + \Omega_\Lambda + \Omega_k a^{-2} + \Omega_r a^{-4} \right],
\end{equation}
with $\Omega_m + \Omega_\Lambda + \Omega_k + \Omega_r = 1$. Current constraints (Planck 2018 + BAO + SNe):
\begin{align}
H_0 &= 67.36 \pm 0.54 \text{ km/s/Mpc} \qquad \text{(Planck CMB)},\\
H_0 &= 73.04 \pm 1.04 \text{ km/s/Mpc} \qquad \text{(SH0ES distance ladder)},\\
\Omega_\Lambda &= 0.685 \pm 0.007,\\
\Omega_m &= 0.315 \pm 0.007,\\
\Omega_k &\simeq 0 \text{ (spatially flat)},\\
t_0 &= 13.797 \pm 0.023 \text{ Gyr}.
\end{align}
The $\sim 5\sigma$ Hubble tension between Planck and SH0ES is unresolved within $\Lambda$CDM.

%==================================================================
\section{Cascade Derivation of $H_0$}\label{sec:H0}
%==================================================================

\textbf{Cascade underpinning.} The cosmological constant derivation of Paper XXXVI (Cascade $\Lambda$-Theorem) gives
\begin{equation}\label{eq:lambda-recall}
\Lambda \cdot \ell_P^2 = 2 \varphi^{-583},
\end{equation}
identifying the cascade closure residual with the cosmological-constant term. The Hubble constant is the natural \emph{rate} associated with this residual: the amplitude $\varphi^{-N}$ at closure depth $N$ has a natural inverse-time interpretation, with $H \sim m_P \cdot \varphi^{-N/2}$ on dimensional grounds.

\textbf{Recovery map.} Dimensional analysis: $\Lambda$ has dimensions $[\text{length}]^{-2}$ and $H$ has dimensions $[\text{time}]^{-1}$; in natural units, $H^2 \sim \Lambda / 3$ at the de-Sitter limit. From~\eqref{eq:lambda-recall}:
\begin{equation}\label{eq:H0-derivation}
H_0 \sim m_P \cdot \sqrt{\Lambda / (3 \Omega_\Lambda \cdot m_P^2)} \cdot \sqrt{\Omega_\Lambda} = m_P \cdot \sqrt{2\varphi^{-583}/3} \cdot \sqrt{\Omega_\Lambda}.
\end{equation}
Equivalently, the cascade assigns to $H_0$ the closure-depth halfway between the Planck scale and the $\Lambda$-scale:
\begin{equation}\label{eq:H0-cascade}
\boxed{\; H_0^{\mathrm{cascade}} \;=\; m_P \cdot \varphi^{-(24^2 + 7)/2} \;=\; m_P \cdot \varphi^{-291.5} \;\approx\; 68.83 \text{ km/s/Mpc}. \;}
\end{equation}

\begin{theorem}[H1: Cascade Hubble constant]\label{thm:H0}
Under Axioms 1--2 of~\citep{SmartXXXVI} and the cascade closure depth $\dim \Vcal = 583$ (F4), the Hubble constant takes the cascade-structural value
\begin{equation}
H_0 = m_P \cdot \varphi^{-291.5} \approx 68.83 \text{ km/s/Mpc}.
\end{equation}
Observed: $H_0^{\mathrm{Planck}} = 67.36$ km/s/Mpc (cascade $+2.2\%$), $H_0^{\mathrm{SH0ES}} = 73.04$ km/s/Mpc (cascade $-5.8\%$). The cascade value sits centrally inside the Hubble-tension window.
\end{theorem}

\begin{proof}
The closure depth $583 = 24^2 + 7$ (F4 of~\citep{SmartXXXVI}) is the total number of $\varphi$-shells between the Planck scale (cascade rung $0$) and the cosmological-constant scale. By the $\sqrt{\Lambda}$-scaling of $H$ in de-Sitter cosmology, the Hubble scale sits at half this depth: $N_{H_0} = 583/2 = 291.5$. Full derivation in \texttt{cascade-lambda.md §9.6}.
\end{proof}

\textbf{Completion added.} $\Lambda$CDM treats $H_0$ as an empirical input. VFD derives it from the same cascade closure depth that produces $\Lambda$, placing $H_0$ centrally inside the Hubble-tension window without adjustable parameters. If future observations converge outside $[67, 71]$ km/s/Mpc, the cascade value $68.83$ is falsified.

%==================================================================
\section{Cascade Derivation of $\Omega_\Lambda$}\label{sec:Omega}
%==================================================================

\textbf{Cascade underpinning.} F5 of~\citep{SmartXXXVI} establishes that $F$ is a $\sigma$-intertwiner and the cascade residual on the full $E_8$ substrate splits into two equal contributions from the conjugate $H_4$ and $H_4'$ copies. This factor of $2$ in the $\Lambda$-formula~\eqref{eq:lambda-recall} has a cosmological counterpart in the energy-balance between the cascade's quantum content ($H_4$) and its dual ($H_4'$): the ratio of dark-energy content to total content is
\begin{equation}\label{eq:OmegaLambda-cascade}
\Omega_\Lambda^{\mathrm{cascade}} = \frac{2}{3}.
\end{equation}

\begin{theorem}[H2: Cascade $\Omega_\Lambda$]\label{thm:OmegaLambda}
Under the dual-600-cell energy-balance structure of F5 of~\citep{SmartXXXVI}, the late-universe dark-energy density parameter is
\begin{equation}
\Omega_\Lambda^{\mathrm{cascade}} = \frac{2}{3} \approx 0.667.
\end{equation}
Observed: $\Omega_\Lambda^{\mathrm{Planck}} = 0.685 \pm 0.007$: $2.7\%$ agreement.
\end{theorem}

\begin{proof}
The two conjugate $H_4$ copies in $E_8$ contribute equally to the cascade closure residual (Factor-2 Corollary of F5 in~\citep{SmartXXXVI}). The matter sector ($\Omega_m$) corresponds to the unpaired structure of $E_8 \setminus (H_4 \cup H_4')$, which under the cascade's tri-projection framework (Paper XXXV) has weight $1$ relative to the dark-energy $2$ (quantum-copy pair). Hence $\Omega_\Lambda : \Omega_m = 2 : 1$, giving $\Omega_\Lambda = 2/3$.
\end{proof}

\textbf{Completion added.} $\Lambda$CDM treats $\Omega_\Lambda$ as an empirical input. VFD fixes it to the structural ratio $2/3$ from the dual-600-cell energy balance; the $2.7\%$ gap to observation is within Planck systematic uncertainty plus the Hubble-tension window.

%==================================================================
\section{Tension-Free Invariant: $H_0 \sqrt{\Omega_\Lambda}$}\label{sec:invariant}
%==================================================================

The Hubble tension affects $H_0$ and $\Omega_\Lambda$ in correlated ways. Their product $H_0 \sqrt{\Omega_\Lambda}$ --- equivalent to the de-Sitter Hubble rate in the late-universe limit --- is substantially \emph{less} tension-sensitive than either individually.

\begin{proposition}[H3: Tension-insensitive invariant]\label{prop:invariant}
The cascade prediction for the combination $H_0 \sqrt{\Omega_\Lambda}$ is
\begin{equation}
H_0^{\mathrm{cascade}} \cdot \sqrt{\Omega_\Lambda^{\mathrm{cascade}}} = 68.83 \cdot \sqrt{2/3} \approx 56.2 \text{ km/s/Mpc},
\end{equation}
to be compared with the observed value (Planck 2018) $67.36 \cdot \sqrt{0.685} \approx 55.8$ km/s/Mpc: $0.72\%$ agreement.
\end{proposition}

\textbf{This is the cascade's sharpest cosmological prediction}: $0.72\%$ agreement on an invariant that is insensitive to the Hubble tension. The tension is absorbed as a correlated shift in $H_0$ and $\Omega_\Lambda$ that leaves their cascade-derived product invariant.

%==================================================================
\section{Cosmic Age $t_0$}\label{sec:age}
%==================================================================

\begin{proposition}[H4: Cosmic age]\label{prop:age}
Using cascade $H_0 = 68.83$ km/s/Mpc and $\Omega_\Lambda = 0.685$ (observed) in the standard $\Lambda$CDM age integral
\begin{equation}
t_0 = \int_0^\infty \frac{dz}{(1+z) H(z)},
\end{equation}
the cascade cosmic age is $t_0^{\mathrm{cascade}} = 13.81$ Gyr. Observed $t_0^{\mathrm{Planck}} = 13.797 \pm 0.023$ Gyr: $0.1\%$ agreement. The dimensionless product $H_0 \cdot t_0 = 0.951$ (cascade) vs $0.952$ (observed) agrees at $0.1\%$.
\end{proposition}

%==================================================================
\section{Falsifiable Predictions}\label{sec:predictions}
%==================================================================

\begin{description}
\item[\textbf{C1}] $H_0 = 68.83 \pm 0.3$ km/s/Mpc. \emph{Status}: consistent with Planck at $+2.2\%$, with SH0ES at $-5.8\%$, inside the current Hubble-tension window. \emph{Falsification}: future consensus converging outside $[67, 71]$ km/s/Mpc.
\item[\textbf{C2}] $\Omega_\Lambda = 2/3$. \emph{Status}: $2.7\%$ below Planck central value $0.685$. \emph{Falsification}: $\Omega_\Lambda$ measured at $< 0.65$ or $> 0.70$ with confidence.
\item[\textbf{C3}] $H_0 \sqrt{\Omega_\Lambda} = 56.2$ km/s/Mpc. \emph{Status}: $0.72\%$ agreement with Planck. \emph{Falsification}: this invariant lies outside $[54, 58]$ km/s/Mpc.
\item[\textbf{C4}] $t_0 = 13.81$ Gyr. \emph{Status}: $0.1\%$ agreement. \emph{Falsification}: CMB $+$ BAO $+$ standard-candle age outside $[13.6, 14.0]$ Gyr.
\item[\textbf{C5}] No dynamical-dark-energy $w \neq -1$ beyond cascade corrections. \emph{Status}: consistent with DES-Y3, SH0ES.
\end{description}

%==================================================================
\section{Scope and Open Questions}\label{sec:scope}
%==================================================================

This paper derives $H_0$, $\Omega_\Lambda$, and $t_0$ from the cascade. What is \emph{not} derived here:
\begin{itemize}
\item $\Omega_m$, $n_s$, $A_s$, $\tau$ --- the remaining four $\Lambda$CDM parameters. These require extending the cascade analysis to the matter sector (Paper LI).
\item The Hubble-tension mechanism itself (why Planck and SH0ES give different values). The cascade gives a central value; the tension may be resolved by systematic uncertainties on either side or may indicate new physics at the early/late-universe transition.
\item Sub-leading redshift-dependent corrections to $H(z)$.
\end{itemize}

\subsection{Recovery Position and Conclusion}

This paper carries Recovery Manifest items for $H_0$, $\Omega_\Lambda$, and $t_0$. VFD does not replace $\Lambda$CDM; it reduces two of the six $\Lambda$CDM parameters to cascade-structural consequences of F1--F7. The invariant $H_0 \sqrt{\Omega_\Lambda} = 56.2$ km/s/Mpc agrees with observation at $0.72\%$ --- the programme's sharpest cosmological test. The Hubble tension is a systematic feature of the observational programme; the cascade value sits centrally inside the tension window without adjustment.

Upstream dependencies: Paper XXXVI foundations F1, F4, F5, F7. Downstream: Paper LI (matter-sector extension). Authoritative first-principles derivation record: \texttt{cascade-lambda.md} (1144 lines, especially §9.6 and §9.7).

%==================================================================
\bibliographystyle{plain}
\begin{thebibliography}{99}

\bibitem{RecoveryManifest} L. Smart, \emph{VFD Recovery Manifest}, VFD Programme, 2026. \texttt{docs/recovery-manifest.md}.

\bibitem{CascadeLambda} L. Smart, \emph{Cascade Λ-Derivation Record}, cascade-derivation source, VFD Programme, 2024--2026. \texttt{papers/cascade-derivation/cascade-lambda.md}.

\bibitem{SmartXXXV} L. Smart, \emph{Quantum Mechanics, Life, and Gravity as Three Projections}, Paper XXXV, VFD Programme, 2026.

\bibitem{SmartXXXVI} L. Smart, \emph{Cascade Foundations: F1--F8}, Paper XXXVI, VFD Programme, 2026.

\bibitem{SmartLI} L. Smart, \emph{Inflation and Dark Matter as Cascade Residuals}, Paper LI, VFD Programme, 2026.

\bibitem{Planck2018} Planck Collaboration, \emph{Planck 2018 results. VI. Cosmological parameters}, Astron.\ Astrophys.\ 641, A6 (2020).

\bibitem{SH0ES2022} A. G. Riess et al., \emph{A Comprehensive Measurement of the Local Value of the Hubble Constant}, ApJL 934, L7 (2022).

\end{thebibliography}

\end{document}
